% Prova di compilazione dei testi corretti del controllo finale: xelatex correzioni_prova.tex
% (i rimandi restano irrisolti: le etichette stanno nei capitoli dell'allegato).
\input{/home/fede/Scrivania/ed/GitHub/Movi-Hackthon/.sciame/20261114-1301-percezione-challenge/f13-allegato/30-finale/preambolo}
\begin{document}
\section*{R1}
Da scheda, con l'avviso che resta dopo il riaggancio: soggetto nascosto entro $3\unita{m}$ senza avviso in 0 fotogrammi su 4324 a 424 pixel di larghezza e in 0 su 4325 a 848, e in 0 su 3608 per risoluzione in una replica; con la configurazione di riferimento, in 180 e 271 degli stessi fotogrammi.

Con l'avviso che resta dopo il riaggancio, un soggetto nascosto entro $3\unita{m}$ non resta senza avviso (0 fotogrammi su 4324 a 424 pixel di larghezza e 0 su 4325 a 848, e 0 su 3608 per risoluzione in una replica); con la configurazione di riferimento restava senza avviso in 180 e 271 di quei fotogrammi (\S\ref{sec:val}).
\section*{R2}
(5) sul suolo vuoto, in $30\unita{s}$, non devono nascere gruppi né tracce confermate; la prova si esegue nella configurazione che si vuole attivare, cioè con i gruppi a celle forti (\S\ref{sec:sti-mobile}): con i gruppi per componenti connesse della configurazione di riferimento, su profondità simulata dalle schede tecniche, sul suolo vuoto nascono 137 e 206 tracce confermate in 9 sequenze (a 424 e 848 pixel di larghezza), e la prova non sarebbe superata \prova{S}. Se cade una delle prime quattro prove il sistema non parte; se cade la quinta restano spenti i gruppi a celle forti, le condizioni aggiuntive del riaggancio, l'avviso senza traccia e l'identità al ritorno, e l'uscita lo dichiara.
\section*{R3}
La frase per settori è una regola di progetto: non è stata generata in alcuna prova, e l'ordine fra i primi due livelli non è stato provato \prova{N}.
\section*{R4, R6, R13, R14 (righe del cap. 1); R5, R11 (righe del cap. 6)}
\begin{table}[H]\centering\footnotesize
\begin{tabularx}{\textwidth}{@{}>{\RaggedRight\arraybackslash}p{2.8cm}>{\RaggedRight\arraybackslash}p{4.1cm}L>{\RaggedRight\arraybackslash}p{1.2cm}>{\RaggedRight\arraybackslash}p{1.05cm}@{}}
\toprule
Risultato & Condizioni & Valore & Prova & Dove\\
\midrule
Osservato distinto dal previsto & da scheda; configurazione di riferimento / gruppi a celle forti & uscita «osservato» con il soggetto nascosto in 8 passaggi su 90 in una prima serie e in 11 e 6 su 140 in una seconda (a 424 e 848 pixel di larghezza) / in 1 e 1 su 140: la soglia 0 non è rispettata & \prova{S} & \S\ref{sec:sti}\\
Regola «mobile» & ideale, $0{,}5\unita{m/s}$ da $2\unita{m}$, $N=2000$; da scheda, pannello fermo, 90 passaggi & «in movimento» entro $1\unita{s}$: 100\,\% ($\ell^-=99{,}81\,\%$); pannello fermo dichiarato mobile in 79 su 90 $[79{,}4;\,93{,}0\,\%]$; in 5 su 90 con la soglia minima di punti portata da 3 a 40, non adottata perché una scatola alta $15\unita{cm}$ a $2\unita{m}$ è allora vista in 3 prove su 10 anziché in 10 su 10; con i gruppi a celle forti in 0 su 140 nel dominio dell'identità al ritorno e in 7 su 10 fuori da esso & \prova{S} & \S\ref{sec:sti}\\
\quad su profondità simulata & da scheda; configurazione di riferimento / gruppi a celle forti e condizioni aggiuntive del riaggancio, soggetto a $0{,}5\unita{m}$ dietro il pannello & 0 su 90 $[0;\,4{,}1\,\%]$ / 176 su 180 $[94{,}4;\,99{,}1\,\%]$ a 424 pixel di larghezza e 174 su 180 $[92{,}9;\,98{,}5\,\%]$ a 848: un tasso in un dominio ristretto, non il superamento di una soglia per ciascun comportamento. Un secondo soggetto nella stessa ombra riceve l'identità del primo (19 casi su 20 da scheda; 2400 su 2400 con modello ideale). L'avviso resta per costruzione; da scheda, in una scena con due persone, 0 fotogrammi su 336 con il primo soggetto fuori dall'avviso; con modello ideale la verifica non è stata eseguita & \prova{S}\prova{N} & \S\ref{sec:val}\\
\bottomrule
\end{tabularx}
\end{table}
\begin{table}[H]\centering\footnotesize
\begin{tabularx}{\textwidth}{@{}>{\RaggedRight}p{4.2cm}L>{\RaggedRight}p{2.1cm}>{\RaggedRight}p{1.9cm}>{\RaggedRight}p{2.3cm}l@{}}
\toprule
Misura & Condizioni & Valore & $N$ & Intervallo 95\,\% & Dato\\
\midrule
Ostacolo fermo entro 15 cm dalla posizione vera & cassa alta $1\unita{m}$ a $1{,}68\unita{m}$; scatola di 15 cm a $2{,}0\unita{m}$ & 84; 10 & 84; 10 fotogrammi (una sequenza ciascuno) & — & II \prova{S}\\
idem, prove d'insieme & corpo esteso e soggetto puntiforme; 12 configurazioni $\times$ 5 comportamenti & 0 & 23\,402 passaggi (6000 soggetti distinti) & 0–0,064\,\% (su 6000) & I\\
\bottomrule
\end{tabularx}
\end{table}
\section*{R7}
\begin{enumerate}
\item \textbf{Ostacolo che nasconde per intero.} Per una persona alta $1{,}70\unita{m}$ a $2{,}3\unita{m}$ dal sensore servono da $1{,}57$ a $1{,}61\unita{m}$ d'ostacolo con il pannello fra $1{,}7$ e $1{,}9\unita{m}$ ($1{,}59\unita{m}$ a $1{,}8\unita{m}$), e $1{,}65\unita{m}$ per una statura di $1{,}75\unita{m}$ con il pannello a $1{,}9\unita{m}$ (Tabella~\ref{tab:geo-amin}).
\end{enumerate}
Dietro un ostacolo più basso dell'altezza minima della Tabella~\ref{tab:geo-amin} (circa $1{,}65\unita{m}$ per una statura di $1{,}75\unita{m}$ con il pannello a $1{,}9\unita{m}$) il sistema tratta come nascosta una persona ancora in parte visibile
\section*{R8}
(a) A sensore fermo del profilo conta solo $T_r$. Con $T_r=1{,}0\unita{s}$ la velocità coperta dall'avviso tempestivo è $1{,}44\unita{m/s}$: una persona a passo normale ($1{,}4\unita{m/s}$) vi rientra con $6\unita{cm}$ di margine sulla distanza, il passo svelto no. È una grandezza diversa dalla velocità coperta dalla fascia di guardia ($1{,}13\unita{m/s}$, condizione~5 del \S\ref{sec:lim-validita}), che il passo normale supera: a $1{,}4\unita{m/s}$ l'avviso è tempestivo per calcolo, ma il contenimento nella fascia non è dimostrato (\S\ref{sec:occ}).
\section*{R9}
\begin{itemize}
\item La fascia di guardia $w$ contiene il centro del soggetto se vale la condizione~\eqref{eq:occ-fascia}, l'ostacolo è fermo e il soggetto resta nel campo e in portata; con $w=0{,}40\unita{m}$ la velocità coperta, per calcolo, è $1{,}13\unita{m/s}$, e $0{,}56\unita{m/s}$ se il bordo dell'ombra sbaglia di $0{,}20\unita{m}$; per contenere l'intero corpo, di raggio $r_c$, la condizione chiede $r_c$ in più e la velocità coperta è minore (Proposizione~\ref{prop:occ-fascia}) \prova{C}.
\end{itemize}
\section*{R14, R15}
Le simulazioni sono di due tipi: \emph{ideale}, cioè con modello ideale del sensore (posizioni al suolo con rumore gaussiano di $5\unita{cm}$, probabilità di rilevamento per fotogramma $P_D$, sensore e ostacolo fermi, un soggetto per ombra), e \emph{da scheda}, cioè su profondità simulata dalle schede tecniche di tre camere a stereo attivo a $424\times240$ pixel e, dove indicato, a $848\times480$, a posa nota. Gli intervalli fra parentesi quadre sono di Wilson al 95\,\%; gli estremi superiori degli intervalli per zero eventi sono arrotondati per eccesso ($2{,}85\,\%$ è scritto $2{,}9\,\%$ e $0{,}0155\,\%$ è scritto $0{,}016\,\%$), gli altri estremi al valore più vicino.

I gruppi a celle forti, le condizioni aggiuntive del riaggancio e l'avviso per le tracce chiuse vicino a un'ombra non fanno parte della configurazione di riferimento: si attivano all'avvio solo se le prove d'ingresso del \S\ref{sec:lim-salute} sono superate, e l'uscita dichiara quali sono attive.
\section*{R10}
Rispetto alla soglia della Tabella~\ref{tab:intro-requisiti} (0 passaggi con un'interruzione), lo zero vale fra i passaggi con l'ingresso riconosciuto ($22\,474$ su $23\,402$; $23\,378$ con la regola aggiuntiva): negli altri non c'è regione, e su tutti i passaggi la soglia non è raggiunta.
\section*{R12}
Si descrive un prototipo di percezione che rileva elementi mobili e ostacoli con una camera di profondità a stereo attivo, frontale e ferma, e che continua a dare un'indicazione dichiarata quando un ostacolo nasconde il soggetto. Il sistema tiene distinte tre informazioni: l'osservato; il non osservato ma previsto, con una regione e un livello; il non noto. Lo stato «libero» è scritto solo per il suolo osservato nel fotogramma corrente; l'avviso di possibile presenza non termina senza osservazione diretta. Nelle ipotesi dichiarate sono dimostrati che ciò che non è osservato non è mai «libero», il contenimento del soggetto nascosto nell'insieme di possibile presenza e la condizione sulla fascia di guardia. In simulazione, con modello ideale del sensore e un soggetto alla volta, la regione di previsione è emessa a ogni ciclo dell'occlusione una volta riconosciuto l'ingresso nell'ombra (0 interruzioni su $23\,402$ passaggi; l'ingresso non è riconosciuto in 928 di essi, in 24 con una regola aggiuntiva, attiva solo dopo una prova d'ingresso); con un ostacolo fermo largo $2\unita{m}$, il bordo dell'ombra noto e il soggetto sul proprio percorso, la regione contiene il soggetto in almeno 9 casi su 10 fino a $2\unita{s}$ e in almeno 85 su 100 fino a $3\unita{s}$; con un ostacolo largo $1\unita{m}$ la soglia di 9 su 10 non è raggiunta, e su profondità simulata dalle schede tecniche questa regola non è stata eseguita. Con un ostacolo fermo l'avviso resta acceso, e nella realizzazione provata resta acceso anche dopo la rimozione dell'ostacolo.
\section*{R17, R18}
valore di progetto, non misurato; va fissato con $60\unita{s}$ di conteggio a vuoto e sotto carico

4 periodi a $15\unita{Hz}$: valore usato nei conti per la conferma «3 misure su 5»; il legame fra i due numeri non è derivato in questo documento

intervallo utile della camera; non misurato

con lo stato iniziale mescolato il rapporto torna a $0{,}94$. Le tavole dicono quali cifre vengono da semi indipendenti; non sono stati ripetuti con lo stato iniziale mescolato 71\,671 dei 175\,802 istanti del controllo sull'area d'avviso

Sistemi industriali di sorveglianza tridimensionale dichiarano un principio analogo (\S\ref{sec:cos-correlati}).
\end{document}
